\subsection{Definition of Vector Space}

\begin{exercise}[1b1]
    Prove that $-(-v) = v$ for every $v \in V$.
\end{exercise}

\begin{sol}
    Let $v \in V$. By the definition of additive inverse, we have
    \[v + (-v) = 0 = -(-v) + (-v).\]
    Since additive inverses are unique, we conclude that $-(-v) = v$.
\end{sol}

\begin{exercise}[1b2]
    Suppose $a \in \f, v \in V$, and $av = 0$. Prove that $a = 0$ or $v = 0$.
\end{exercise}

\begin{sol}
    If $a = 0$, then we are done. Suppose $a \neq 0$. Then
    \[v = 1v = \frac1a(av) = \frac1a \cdot 0 = 0. \qh\]
\end{sol}

\begin{exercise}[1b3]
    Suppose $v, w \in V$. Explain why there exists a unique $x \in V$ such that $v + 3x = w$.

    \remark{The solution says ``dividing by 3'' in the sense that we're dividing by the scalar 3 in the field $\f$, not a vector in $V$.}
\end{exercise}

\begin{sol}
    Solving for $x$ explicitly, we obtain
    \[3x = w - v \implies x = \frac13(w - v).\]
    This shows that there exists $x \in V$ such that $v + 3x = w$. To see that $x$ is unique, suppose there exists another $y \in V$ such that $v + 3y = w$. Then dividing by 3, we get
    \[3x = w - v = 3y \implies x = y. \qh\]
\end{sol}

\begin{exercise}[1b4]
    The empty set is not a vector space. The empty set fails to satisfy only one of the requirements listed in the definition of a vector space. Which one?
\end{exercise}

\begin{sol}
    The empty set fails to satisfy the requirement that there exists a zero vector $0 \in V$ such that $v + 0 = v$ for all $v \in V$ since it contains no elements.
\end{sol}

\begin{exercise}[1b5]
    Show that in the definition of a vector space, the additive inverse condition can be replaced with the condition that
    \[0v = 0 \text{ for all } v \in V.\]
    Here, the 0 on the left side is the number 0, and the 0 on the right side is the additive identity in $V$.
\end{exercise}

\begin{sol}
    Suppose that $0v = 0$ for all $v \in V$. Let $v \in V$. Then
    \[v + (-1)v = 1v + (-1)v = (1 + (-1))v = 0v = 0.\]
    Hence, $(-1)v$ is the additive inverse of $v$.
    
    Conversely, suppose that for every $v \in V$, there exists an additive inverse $-v \in V$ such that $v + (-v) = 0$. Then $0v = (0 + 0)v = 0v + 0v$. Adding $-0v$ to both sides, we get $0 = 0v$. Hence, $0v = 0$ for all $v \in V$.
\end{sol}

\begin{exercise}[1b6]
    Let $\infty$ and $-\infty$ denote two distinct objects, neither of which is in $\r$. Define an addition and scalar multiplication on $\r \cup \set{\infty, -\infty}$ as you could guess from the notation. Specifically, the sum and product of two real numbers is as usual, and for $t \in \r$, define
    \[t\infty = 
    \begin{cases}
        -\infty & \text{if } t < 0, \\
        0 & \text{if } t = 0, \\
        \infty & \text{if } t > 0,
    \end{cases} \quad
    t(-\infty) = 
    \begin{cases}
        \infty & \text{if } t < 0, \\
        0 & \text{if } t = 0, \\
        -\infty & \text{if } t > 0,
    \end{cases}\]
    and
    \begin{align*}
        t + \infty & = \infty + t = \infty + \infty = \infty, \\
        t + (-\infty) & = (-\infty) + t = (-\infty) + (-\infty) = -\infty, \\
        \infty + (-\infty) & = (-\infty) + \infty = 0.
    \end{align*}
    With these operations of addition and scalar multiplication, is $\r \cup \set{\infty, -\infty}$ a vector space over $\r$? Explain.
\end{exercise}

\begin{sol}
    No, $\r \cup \set{\infty, -\infty}$ is not a vector space over $\r$ as it does not satisfy the associative property of addition. Let $t \in \r \setminus \set 0$. Then
    \[(t + \infty) + (-\infty) = \infty + (-\infty) = 0 \neq t = t + 0 = t + (\infty + (-\infty)). \qh\]
\end{sol}

\begin{exercise}[1b7]
    Suppose $S$ is a nonempty set. Let $V^S$ denote the set of functions from $S$ to $V$. Define a natural addition and scalar multiplication on $V^S$, and show that $V^S$ is a vector space with these definitions.
\end{exercise}

\begin{sol}
    Let $f, g \in V^S$ and $t \in \f$. Define addition and scalar multiplication as follows:
    \[(f + g)(s) = f(s) + g(s), \quad (tf)(s) = t(f(s))\]
    for all $s \in S$. We will show that $V^S$ is a vector space with these operations.
    \begin{enumerate}
        \item \textbf{Commutativity}: Let $f, g \in V^S$. Then for all $s \in S$,
        \[(f + g)(s) = f(s) + g(s) = g(s) + f(s) = (g + f)(s).\]
        \item \textbf{Associativity}: Let $f, g, h \in V^S$ and $t, u \in \f$. Then for all $s \in S$,
        \begin{align*}
            ((f + g) + h)(s) &= (f + g)(s) + h(s) \\
            &= (f(s) + g(s)) + h(s) \\
            &= f(s) + (g(s) + h(s)) \\
            &= f(s) + (g + h)(s) \\
            &= (f + (g + h))(s),
        \end{align*}
        and
        \begin{align*}
            ((tu)f)(s) & = (tu)(f(s)) \\
            & = t(u(f(s))) \\
            & = t(uf)(s) \\
            & = (t(uf))(s).
        \end{align*}
        \item \textbf{Additive identity}: Let $0 \in V$ be the additive identity in $V$. Define the zero function $0_S \in V^S$ by $0_S(s) = 0$ for all $s \in S$. Then for all $f \in V^S$ and $s \in S$,
        \[(f + 0_S)(s) = f(s) + 0_S(s) = f(s) + 0 = f(s) = 0 + f(s) = (0_S + f)(s).\]
        Hence, $0_S$ is the additive identity in $V^S$.
        \item \textbf{Additive inverse}: Let $f \in V^S$. Define the function $-f \in V^S$ by $(-f)(s) = -f(s)$ for all $s \in S$. Then for all $s \in S$,
        \[(f + (-f))(s) = f(s) + (-f)(s) = f(s) + (-f(s)) = 0 = 0_S(s).\]
        Hence, $-f$ is the additive inverse of $f$ in $V^S$.
        \item \textbf{Multiplicative identity}: For all $f \in V^S$ and $s \in S$,
        \[(1f)(s) = 1(f(s)) = f(s).\]
        \item \textbf{Distributive properties}: Let $f, g \in V^S$ and $t, u \in \f$. Then for all $s \in S$,
        \begin{align*}
            ((t + u)f)(s) & = (t + u)(f(s)) \\
            & = t(f(s)) + u(f(s)) \\
            & = (tf)(s) + (uf)(s) \\
            & = (tf + uf)(s),
        \end{align*}
        and
        \begin{align*}
            (t(f + g))(s) & = t((f + g)(s)) \\
            & = t(f(s) + g(s)) \\
            & = t(f(s)) + t(g(s)) \\
            & = (tf)(s) + (tg)(s) \\
            & = (tf + tg)(s). \qh
        \end{align*}
    \end{enumerate}
\end{sol}

\begin{exercise}[1b8]
    Suppose $V$ is a real vector space.
    \begin{itemize}
        \item The \textbf{complexification} of $V$, denoted by $V_{\c}$, equals $V \times V$. An element of $V_{\c}$ is an ordered pair $(u, v)$, where $u, v \in V$, but we write this as $u + iv$.
        \item Addition on $V_{\c}$ is defined by
        \[(u_1 + iv_1) + (u_2 + iv_2) = (u_1 + u_2) + i(v_1 + v_2)\]
        for all $u_1, u_2, v_1, v_2 \in V$.
        \item Complex scalar multiplication on $V_{\c}$ is defined by
        \[(a + ib)(u + iv) = (au - bv) + i(av + bu)\]
        for all $a, b \in \r$ and $u, v \in V$.
    \end{itemize}
    Prove that with the definitions of addition and scalar multiplication as above, $V_{\c}$ is a complex vector space.
\end{exercise}

\begin{sol}
    We prove that $V_{\c}$ is a complex vector space by verifying the vector space axioms.
    \begin{enumerate}
        \item \textbf{Commutativity}: Let $u_1 + v_1i, u_2 + v_2i \in V_{\c}$. Then
        \[(u_1 + v_1i) + (u_2 + v_2i) = (u_1 + u_2) + i(v_1 + v_2) = (u_2 + u_1) + i(v_2 + v_1) = (u_2 + v_2i) + (u_1 + v_1i).\]
        \item \textbf{Associativity}: Let $u_1 + v_1i, u_2 + v_2i, u_3 + v_3i \in V_{\c}$. Then
        \begin{align*}
            ((u_1 + v_1i) + (u_2 + v_2i)) + (u_3 + v_3i) & = ((u_1 + u_2) + i(v_1 + v_2)) + (u_3 + v_3i) \\
            & = ((u_1 + u_2) + u_3) + i((v_1 + v_2) + v_3) \\
            & = (u_1 + (u_2 + u_3)) + i(v_1 + (v_2 + v_3)) \\
            & = (u_1 + v_1i) + ((u_2 + v_2i) + (u_3 + v_3i)).
        \end{align*}
        Similarly, let $a + bi, c + di \in \c$ and $u + vi \in V_{\c}$. Then
        \begin{align*}
            ((a + bi)(c + di))(u + vi) & = ((ac - bd) + i(ad + bc))(u + vi) \\
            & = ((ac - bd)u - (ad + bc)v) + i((ac - bd)v + (ad + bc)u) \\
            & = (a(cu - dv) - b(du + cv)) + i(a(du + cv) + b(cu - dv)) \\
            & = (a + bi)((cu - dv) + i(du + cv)) \\
            & = (a + bi)((c + di)(u + vi)).
        \end{align*}
        \item \textbf{Additive identity}: Let $0 \in V$ be the additive identity in $V$. Then for all $u + vi \in V_{\c}$,
        \[(u + vi) + (0 + 0i) = (u + 0) + i(v + 0) = u + vi = (0 + 0i) + (u + vi).\]
        Hence, $0 + 0i$ is the additive identity in $V_{\c}$.
        \item \textbf{Additive inverse}: Let $u + vi \in V_{\c}$. Then
        \[(u + vi) + (-u + (-v)i) = (u - u) + i(v - v) = 0 + 0i.\]
        Hence, $-u + (-v)i$ is the additive inverse of $u + vi$ in $V_{\c}$.
        \item \textbf{Multiplicative identity}: For all $u + vi \in V_{\c}$,
        \[(1 + 0i)(u + vi) = (1u - 0v) + i(1v + 0u) = u + vi.\]
        \item \textbf{Distributive properties}: Let $a + bi, c + di \in \c$ and $u_1 + v_1i, u_2 + v_2i \in V_{\c}$. Then
        \begin{align*}
            (a + bi)((u_1 + v_1i) + (u_2 + v_2i)) & = (a + bi)((u_1 + u_2) + i(v_1 + v_2)) \\
            & = (a(u_1 + u_2) - b(v_1 + v_2)) + i(a(v_1 + v_2) + b(u_1 + u_2)) \\
            & = (au_1 - bv_1) + i(av_1 + bu_1) + (au_2 - bv_2) + i(av_2 + bu_2) \\
            & = (a + bi)(u_1 + v_1i) + (a + bi)(u_2 + v_2i),
        \end{align*}
        and
        \begin{align*}
            ((a + bi) + (c + di))(u_1 + v_1i) & = ((a + c) + (b + d)i)(u_1 + v_1i) \\
            & = ((a + c)u_1 - (b + d)v_1) + i((a + c)v_1 + (b + d)u_1) \\
            & = (au_1 - bv_1) + i(av_1 + bu_1) + (cu_1 - dv_1) + i(cv_1 + du_1) \\
            & = (a + bi)(u_1 + v_1i) + (c + di)(u_1 + v_1i). \qh
        \end{align*}
    \end{enumerate}
\end{sol}